% Settings for drawing the diagrams in TikZ

\newcommand{\tadpoleangle}{130}
\newlength{\tadpolesize}\setlength{\tadpolesize}{1cm}
\newlength{\bubblesize}\setlength{\bubblesize}{.33cm}
\newcommand{\bubbleaspect}{.8}
\newlength{\photonlength}\setlength{\photonlength}{.8cm}
\newcommand{\diagramscale}{1.25}
\newlength{\vertsize}\setlength{\vertsize}{2.25pt}
\tikzset{
    Ultra thick/.style={line width=3pt},
    miter cap/.style={{Triangle Cap[]}-{Triangle Cap[]}},
    pion/.style = { very thick },
    photon/.style = { thick, decorate, decoration={snake,segment length=.26\photonlength, amplitude=+.05\photonlength} },
    notoph/.style = { thick, decorate, decoration={snake,segment length=.26\photonlength, amplitude=-.05\photonlength} },
    diagrcut/.style={red,decorate,decoration={zigzag,segment length=1mm,amplitude=.25mm}},
    %
    bubble2/.style = {xscale=.8, transform shape},
    elbbub2/.style = {xscale=1.25, transform shape},
    bubble3/.style = {xscale=.64, transform shape},
    threeloop/.style = {scale=1.2, transform shape},
    %
    NLO/.pic = {
        \fill (0,0) circle[radius=\vertsize];
    },
    NNLO/.pic = {
        \fill (-\vertsize,-\vertsize) rectangle (\vertsize, \vertsize);
    },
    NNNLO/.pic = {
        \fill (90:2\vertsize) -- (210:2\vertsize) -- (330:2\vertsize) -- cycle;
    },
    %
    pics/tadpole/.style = {
        code = {
            \draw[pion] (0,0) .. controls (\tadpoleangle:\tadpolesize) and (180-\tadpoleangle:\tadpolesize) .. (0,0)
                foreach \i in {#1} { node[pos=.\i] (node\i) {} }
                ;
        }
    },
    pics/bubble/.style = {
        code = {
            \draw[pion] (0,0) ellipse [x radius=\bubblesize, y radius=\bubbleaspect\bubblesize];
            \foreach \i in {#1} {
                \coordinate (node\i) at
                (canvas polar cs:angle=\i, x radius=\bubblesize, y radius=\bubbleaspect\bubblesize) {};
            }
        }
    },
    pics/bubble/.default = {0}, % For some reason tadpole is fine with empty, but bubble needs default
    bubblejoint/.pic = {
        \fill[white] (-\bubblesize,-\bubblesize) rectangle (\bubblesize,\bubblesize);
        \draw[pion] (-\bubblesize,+\bubbleaspect\bubblesize) .. controls (0,+\bubbleaspect\bubblesize) and (0,-\bubbleaspect\bubblesize) .. (+\bubblesize,-\bubbleaspect\bubblesize);
        \draw[pion] (+\bubblesize,+\bubbleaspect\bubblesize) .. controls (0,+\bubbleaspect\bubblesize) and (0,-\bubbleaspect\bubblesize) .. (-\bubblesize,-\bubbleaspect\bubblesize);
    },
}

\newcounter{diagram}
\newenvironment{diagram}[1][]{%
%         \stepcounter{diagram}\arabic{diagram}
        \begin{tikzpicture}[scale=\diagramscale, transform shape, baseline={(0,0)},#1]%
        \path[use as bounding box] (-\photonlength,-\bubblesize) rectangle (+\photonlength,+\bubblesize);
    }{%
        % Smash vertically, keep horizontal extent
        \end{tikzpicture}%
    }

    
% P. Tol's colorblind-friendly plot colors
\definecolor{plotI}{HTML}{0077BB}
\definecolor{plotII}{HTML}{33BBEE}
\definecolor{plotIII}{HTML}{009988}
\definecolor{plotIV}{HTML}{EE7733}
\definecolor{plotV}{HTML}{CC3311}
\definecolor{plotVI}{HTML}{DDAA33}
\definecolor{plotVII}{HTML}{AA3377}

% Colors for various kinematic regions
\colorlet{spacelike}{plotI}
\colorlet{subthr}{plotIV}
\colorlet{twopi}{plotIII}
\colorlet{fourpi}{plotVII}
\colorlet{shaded}{gray!30}
